%---------------------------------------------------------------------
%
%                        frontmatter/resumen.tex
%
%---------------------------------------------------------------------
%
% El resumen en castellano y sus palabras clave. Es el UNICO texto en
% castellano de la memoria: V.4 permite escribirla integramente en ingles,
% pero V.2b exige el resumen y un maximo de 10 palabras clave en los dos
% idiomas. main.tex lo envuelve en otherlanguage{spanish} para que babel
% lo silabee como castellano.
%
% Media pagina. Debe incluir el titulo en castellano.
%
%---------------------------------------------------------------------

\chapter*{Resumen}
\addcontentsline{toc}{chapter}{Resumen}

\section*{\tituloPortadaVal}

\todo{Escribir el resumen, de media pagina, cuando esten los resultados.
Debe decir, en este orden: (1) que dos canciones "se parezcan" no tiene
una definicion operativa y la literatura usa el genero como sustituto;
(2) que este trabajo construye un grafo de similitud a partir de
representaciones del audio y lo evalua en dos niveles, sondas sobre la
representacion y calidad del grafo; (3) que modelos y que datos se usan
(MFCC, CNN, redes siamesa y de tripletes, MERT; FMA small y
MagnaTagATune); (4) el resultado principal, incluyendo el grado de
acuerdo con los juicios humanos de similitud; y (5) la conclusion.}

\section*{Palabras clave}

\noindent \todo{Maximo 10 palabras clave separadas por comas. Propuesta:
recuperacion de informacion musical, similitud musical, embeddings de
audio, grafos de similitud, vecinos mas cercanos, evaluacion, MFCC,
aprendizaje de metricas, juicios humanos de similitud.}
